\subsection{补算：欧姆定律中的四个符号和单位}
重新把电子方程按代数而非记忆整理：
\begin{align}
 en(\E+\vel_e\times\B)
 &=-m_enD_e\vel_e-\grad\cdot\mathsf P_e+\vct R_e,\\
 \E+\vel_i\times\B
 &=(\vel_i-\vel_e)\times\B
   -{\grad\cdot\mathsf P_e\over en}-{m_e\over e}D_e\vel_e
   +{\vct R_e\over en}.
\end{align}
压力项和惯性项的负号来自电子负电荷；Hall项的正号来自左端把电子速度换成离子速度。对经典动量阻力，$\vct R_e$指向离子相对电子的运动方向，因此$R_e/(en)=\eta\J$是正号。

四项的单位都要等于电场。例如
\begin{equation}
 \left[{\grad p_e\over en}\right]
 ={{\rm N\,m^{-3}}\over{\rm C\,m^{-3}}}
 ={\rm N\,C^{-1}}={\rm V\,m^{-1}},\qquad
 \left[{m_e\over e}D_e\vel_e\right]
 ={{\rm kg}\over{\rm C}}{{\rm m}\over{\rm s^2}}.
\end{equation}
不能把某个代码中的无量纲碰撞频率直接代入SI电阻率$\eta$的位置。

\subsection{把密度变化也留在电子惯性中}
式\eqref{eq:inertiaexpanded}还可以进一步展开。令$D_i=\partial_t+\vel_i\cdot\grad$，用离子连续性$D_in=-n\theta_i$：
\begin{align}
 D_i\vct w&=D_i\left({\J\over en}\right)
 ={D_i\J\over en}-{\J\over en^2}D_in
 ={D_i\J+\J\theta_i\over en},\\
 (\vct w\cdot\grad)\vct w
 &={(\J\cdot\grad)\J\over e^2n^2}
 -{\J(\J\cdot\grad n)\over e^2n^3}.
\end{align}
所以电子惯性电场完整写成
\begin{align}
 -{m_e\over e}D_e\vel_e
 ={}&-{m_e\over e}D_i\vel_i
 +{m_e\over e^2n}(D_i\J+\J\theta_i)
 +{m_e\over e^2n}(\J\cdot\grad)\vel_i\\
 &-{m_e\over e^3n^2}(\J\cdot\grad)\J
 +{m_e\over e^3n^3}\J(\J\cdot\grad n).\nonumber
\end{align}
最后两项来自电子相对离子的对流。只保留$\partial_t\J$，相当于同时舍弃流体加速度、压缩、剪切、自对流和密度梯度；只有在明确的线性化或局部排序下才可能合理。常密度只能删去含$\grad n,\theta_i$的相关项，不能自动删去所有空间对流。

\subsection{为什么小电子质量可以乘上一个重要项}
电子层中若$u_e\sim J/(en)$、$J\sim B/(\mu_0\delta)$，则
\begin{equation}
 u_e\sim {B\over\mu_0en\delta}.
\end{equation}
在具有相应非零分量的区域，以$u_e^2/\delta$估算惯性、以$u_eB$估算对流电场，得
\begin{equation}
 { (m_e/e)u_e^2/\delta\over u_eB}
 \sim{m_e\over\mu_0ne^2\delta^2}
 ={d_e^2\over\delta^2}.
\end{equation}
因此$\delta\sim d_e$时电子惯性可以为主阶，尽管$m_e/m_i$很小。这里不能在X点对两个都为零的特定分量作比值；X点要直接使用式\eqref{eq:xpoint}检查留下的项。

\subsection{手算两个X点例子}
\paragraph{例1：标量压力给不出二维重联电场。}
若$P_{e,ij}=p_e\delta_{ij}$且$\partial_z=0$，则
$P_{exz}=P_{eyz}=0$，$\partial_zP_{ezz}=0$。
在面内磁零点，$(\vel_e\times\B)_z=u_{ex}B_y-u_{ey}B_x=0$。一个很强的$B_z$不会出现在这个分量中。若再令惯性和电阻为零，$E_z(X)$便只能为零。

\paragraph{例2：非对角压力如何支撑有限电场。}
取一个局地光滑的压力张量，使X点附近
$P_{exz}=a x$、$P_{eyz}=b y$，其余分量足够大以保持正定。若该处惯性和电阻贡献可忽略，
\begin{equation}
 E_z(X)=-{a+b\over en}.
\end{equation}
压力在X点本身的两个非对角值都是零，梯度却可以非零。这解释了为什么只看$P_{xz}(X)$的数值还不够，必须计算其空间导数。例子只是局部支撑项演示，不是已经求解了完整全局重联平衡。
